\chapter{Gramáticas libres de contexto múltiple paralelas}

\section{Definiciones}

Las gramáticas libres de contexto múltiples (MCFG) y libres de contexto múltiples paralelas (pMCFG) son generalizaciones de las gramáticas libres de contexto (CFG) que permiten la generación de una familia de lenguajes más amplia que estas, pero más reducida que la de los lenguajes dependientes del contexto \cite{Seki1991OnMC}. La idea central es extender las CFG de forma tal que los símbolos no terminales puedan generar tuplas de cadenas en lugar de una sola cadena \cite{Seki1991OnMC}.

Estas gramáticas contienen reglas de producción capaces de especificar cómo combinar tuplas de cadenas para generar nuevas tuplas, mediante la aplicación de una función intermedia que reordena y concatena las cadenas generadas por los símbolos no terminales. A continuación se definen dichas funciones, las cuales serán denominadas \lquoti funciones de reordenamiento\rquoti \space \cite{ParsingBeyCFG}:

\begin{definition}
    Una función de reordenamiento es una función que recibe tuplas de cadenas y devuelve una nueva tupla de cadenas resultantes a partir de la concatenación de las cadenas recibidas.

    Formalmente, $f$ es una función de reordenamiento sobre un alfabeto $T$ si existe un número entero $k$ y $d_i > 0$ para $0 \leq i \leq k$ tal que $f: (T^*)^{d_1} \times  \ldots \times (T^*)^{d_k} \to (T^*)^{d_0}$ tal que las componentes de $f(x_1, \ldots, x_k)$ son concatenaciones de un número finito de símbolos de $T$ y de las componentes $y_{ij}$ de las tuplas $x_i$ ($1 \leq i \leq k$, $1 \leq j \leq d_i$).
\end{definition}

Por ejemplo, dado el alfabeto $T = \{a, b\}$, las funciones $f: T^* \to (T^*)^2$, $g: T^* \times T^* \to T^*$ y $h: (T^*)^2 \times T^* \to (T^*)^2$ definidas como:

$$
    f(x_1) = (x_1x_1, ax_1) \\
    \quad g(x_1, x_2) = x_2ax_1b \\
    \quad h((x_{11}, x_{12}), x_{2}) = (x_{11}x_{2}, x_{12})
$$
% cualquier comentario sobre los ejemplos anteriores y la forma en la que escribí los argumentos es bienvenido

son funciones de reordenamiento. En los casos particulares de $g$ y $h$ se tiene que las componentes $x_{ij}$ de las tuplas $x_i$ aparecen a lo sumo una vez en las componentes de la tupla resultante de aplicar $g$ y $h$. A las funciones de reordenamiento que cumplen esta condición se les llama funciones de reordenamiento lineales. La función $f$ del ejemplo anterior no es lineal, mientras que las funciones $g$ y $h$ sí lo son.

Las funciones de reordenamiento permiten combinar las cadenas generadas por distintos símbolos no terminales. A partir de estas se definen las gramáticas libres de contexto múltiples paralelas.

\begin{definition}
    Una gramática libre de contexto múltiple paralela (pMCFG) es una tupla $G = (N, T, F, P, S)$ donde:

    \begin{itemize}
        \item $N$ es un conjunto finito de símbolos llamados símbolos no terminales, cada $A \in N$ tiene una dimensión asociada $dim(A) \geq 1$.
        \item $T$ es un conjunto finito de símbolos llamados terminales. Se cumple que $N \cap T = \emptyset$.
        \item $F$ es un conjunto finito de funciones de reordenamiento.
        \item $P$ es un conjunto finito de reglas de la forma $A_0 \to f(A_1, \ldots, A_k)$ con $k \geq 0, f \in F$ tal que $f: (T^*)^{dim(A_1)} \times \ldots \times (T^*)^{dim(A_k)} \to (T^*)^{dim(A_0)}$.
        \item $S$ es un símbolo no terminal inicial con $dim(S) = 1$.
    \end{itemize}
\end{definition}

% la razón por la cual no defino la derivación como tal, es que en ningún lado la encuentro como tal, en plan A => algo1 => algo2 

El proceso de derivación en una pMCFG es similar al de una CFG, pero con la aplicación de las funciones de reordenamiento para combinar las tuplas de cadenas generadas por los símbolos no terminales. La pMCFG genera un lenguaje que es el conjunto de todas las cadenas terminales que pueden derivarse a partir del símbolo inicial $S$.


\begin{definition}
    Sea $G = (N, T, F, P, S)$ una pMCFG. Se denota por $L(A)$ con $A \in N$ al conjunto de tuplas de cadenas de símbolos terminales que pueden ser derivadas a partir del símbolo no terminal $A$.

    Para todo $A \in N$, $L(A)$ se define de la siguiente manera:
    \begin{itemize}
        \item Para toda regla $A \to f() \in P$, se tiene que $f() \in L(A)$.
        \item Para toda regla $A \to f(A_1, \ldots, A_k) \in P$ con $k \geq 1$, y todas las tuplas de cadenas $x_i \in L(A_i)$ para $1 \leq i \leq k$, se tiene que $f(x_1, \ldots, x_k) \in L(A)$.
        \item No hay otras tuplas de cadenas en $L(A)$.
    \end{itemize}
    El lenguaje generado por la pMCFG $G$ es el conjunto de cadenas terminales que pueden ser derivadas a partir del símbolo inicial $S$, es decir, $L(G) = \{ w \in T^* \suchthat (w) \in L(S) \}$.
\end{definition}

Por ejemplo, $G = (N, T, F, P, S)$ con $N = \{S\}, T = \{a\}, F = \{f, g\}$ y $P = \{S \to f(), S\to g(S)\}$ donde $f() = a$ y $g(x) = xx$ es una pMCFG que genera el lenguaje $\{a^{2^n} \suchthat n \geq 0\}$ \cite{Seki1991OnMC}.

% necesito clases de como escribir la matemática inline y que no haga word-break al final del renglón xd

En \cite{Seki1991OnMC} se definen también las gramáticas libres de contexto múltiples (MCFG):

\begin{definition}
    Una gramática libre de contexto múltiple (MCFG) es una tupla $G = (N, T, F, P, S)$ donde $N, T, F, P, S$ cumplen las mismas condiciones que en la definición de pMCFG, pero con la restricción adicional de que todas las funciones de reordenamiento en $F$ son funciones de reordenamiento lineales.
\end{definition}

El lenguaje generado por una MCFG se define de la misma manera que el lenguaje generado por una pMCFG.

Por ejemplo, $G = (N, T, F, P, S)$ con $N = \{S, A\}, T = {0, 1}, F=\{f_1, f_2, f_3, f_4, f_5\}$ y:

$$
    P = \{
    S \to f_1(A), \quad
    A \to f_2(A), \quad
    A \to f_3(A), \quad
    A \to f_4(), \quad
    A \to f_5()
    \}
$$

con funciones definidas como:
$$
    f_1[(x, y, z)] = (xyz) \quad
    f_2[(x, y, z)] = (ax, ay, az) \quad
    f_3[(x, y, z)] = (bx, by, bz)
$$
$$
    f_4[] = (a, a, a) \quad
    f_5[] = (b, b, b)
$$

es una MCFG que genera el lenguaje $\{www \suchthat w \in \{a, b\}^+\}$ \cite{ParsingBeyCFG}

En siguientes secciones se estudia cómo resolver el problema del vacío de las pMCFG y la intersección de estas con lenguajes regulares y su importancia para este trabajo.

\section{Problema del vacío para pMCFG y MCFG}

Aunque la familia de lenguajes que generan las pMCFG es estrictamente más amplia que la de las CFG \cite{Seki1991OnMC}, resolver el problema del vacío para una pMCFG es igual de complejo que resolver el mismo problema para una CFG. De hecho, el algoritmo para resolver el problema del vacío para una pMCFG es similar al algoritmo análogo para una CFG expuesto en \cite{automataTheory}.

Para exponer dicho algoritmo es necesario definir la noción de símbolo no terminal generador en una pMCFG.

\begin{definition}
    Sea $G = (N, T, F, P, S)$ una pMCFG. Un símbolo no terminal $A \in N$ es generador si y solo si $L(A) \neq \emptyset$.
\end{definition}

Como consecuencia directa, para una pMCFG $G$ con símbolo inicial $S$ se tiene que $L(G) \neq \emptyset$ si y solo si $S$  es generador. Esto hace posible que determinar si $L(G)$ es vacío, sea equivalente a computar el conjunto de de no terminales generadores de $G$ y comprobar si $S$ pertenece o no a este conjunto.

El algoritmo para determinar si $L(G)$ es vacío consiste en construir iterativamente el conjunto de símbolos generadores. Inicialmente, se consideran generadores todos los símbolos no terminales que poseen una regla de producción de la forma $A \to f()$. Luego, mientras existan reglas cuyos símbolos no terminales del lado derecho ya hayan sido identificados como generadores, se añade el símbolo del lado izquierdo al conjunto.

Este procedimiento replica directamente la definición del conjunto $L(A)$, ya que un símbolo no terminal es generador si y solo si puede obtenerse a partir de reglas cuyos argumentos generan tuplas terminales.

\begin{algorithm}[H]
    \caption{Determinación de símbolos generadores en una pMCFG}
    \KwIn{Una pMCFG $M = (N, T, F, P, S)$}
    \KwOut{Conjunto de símbolos generadores} $Gen \gets \{ A \in N \suchthat A \to f() \in P \}$\; $cambio \gets$ verdadero\;
    \While{$cambio$}{ $cambio \gets$ falso\;
        \ForEach{regla $A \to f(A_1, \ldots, A_k) \in P$}{
            \If{$A_1, \ldots, A_k \in Gen$ \textbf{y} $A \notin Gen$}{
                añadir $A$ a $Gen$\;
                $cambio \gets$ verdadero\;
            }
        }
    }
    \Return $Gen$\;
\end{algorithm}

La correctitud del algoritmo es directa dada la definición de $L(A)$. En efecto, cada símbolo añadido al conjunto corresponde a un símbolo no terminal capaz de generar al menos una tupla terminal, y recíprocamente, todo símbolo generador será eventualmente añadido por el algoritmo. Por lo tanto, al finalizar la ejecución del algoritmo, el conjunto calculado coincide exactamente con el conjunto de símbolos generadores de la gramática. En consecuencia, $ L(G) \neq \emptyset \iff S \in Gen$.

Para analizar la complejidad del algoritmo, considérese el tamaño de una pMCFG $G$ definido como el número total de ocurrencias de símbolos no terminales en las reglas de producción de $G$.

En cada iteración del ciclo, el algoritmo recorre todas las reglas de producción y, para cada regla $A \to f(A_1, \ldots, A_k)$ verifica si todos los símbolos $A_1, \cdots, A_k$ pertenecen a $Gen$. Por lo tanto, el costo total de cada iteración es lineal en el tamaño de la gramática.

Además, en cada iteración se añade al menos un símbolo no terminal nuevo a $Gen$, o el algoritmo termina. Como existen a lo sumo $|N|$ símbolos no terminales, el número total de iteraciones es acotado por $|N|$.

En consecuencia, el algoritmo tiene complejidad temporal $O(|N| \cdot |G|)$, y, dado que $|N| \leq |G|$, también puede expresarse como $O(|G|^2)$.

Aunque el algoritmo presentado tiene complejidad polinomial, este puede optimizarse de manera análoga al algoritmo clásico para determinar los símbolos generadores en una CFG \cite{automataTheory}. En lugar de recorrer repetidamente todas las reglas de producción, es posible mantener, para cada regla, la cantidad de símbolos no terminales de su lado derecho que aún no han sido identificados como generadores. Cada vez que un nuevo símbolo se añade a $Gen$, se actualizan únicamente las reglas que dependen de este símbolo. De esta forma, cada regla y cada ocurrencia de un símbolo no terminal son procesadas un número constante de veces. Por lo tanto, el problema del vacío para una pMCFG puede resolverse en tiempo lineal respecto al tamaño de la gramática \cite{Seki1991OnMC}.

\section{Intersección de pMCFG y MCFG con lenguajes regulares}

% añadir ref 
En el capítulo 2 se demuestra que es posible construir el lenguaje de las interpretaciones que satisfacen una fórmula intersectando un autómata finito determinista con el lenguaje $L_{0, 1, d}$. Esto permite que sea posible resolver el problema de la satisfacibilidad booleana decidiendo si el lenguaje resultante de dicha intersección es vacío o no. Dado que el problema del vacío para una pMCFG es decidible en tiempo polinomial, es de interés estudiar la intersección de una pMCFG o una MCFG con un lenguaje regular, y analizar si el resultado de esa intersección pertenece a la misma familia de lenguajes.

\subsection{Intersección de MCFG con DFA}

En \cite{Seki1991OnMC} se demuestra que la intersección de una MCFG con un lenguaje regular es una MCFG. La idea central consiste en enriquecer cada símbolo no terminal $A$ de la MCFG con información sobre los estados recorridos por un DFA al reconocer las componentes de las tuplas generadas por $A$. Si $A$ tiene dimensión $d$, entonces los nuevos símbolos no terminales tienen la forma
$$
    A[(p_1, q_1), \ldots, (p_d, q_d)]
$$

donde cada par $(p_i,q_i)$ representa los estados inicial y final del recorrido del DFA sobre la $i$-ésima componente de la tupla.

Intuitivamente, el símbolo $A[(p_1, q_1), \ldots, (p_d, q_d)]$ en la nueva gramática genera exactamente las mismas tuplas que $A$, pero restringiendo cada componente a cadenas que llevan al DFA desde el estado $p_i$ al estado $q_i$.

% TODO: hablar de las reglas de producción y la condición de conexión

Las reglas de producción de la nueva gramática se construyen a partir de las reglas de producción de la MCFG original. Considérese una regla de producción
$$
    A_0 \to f(A_1, \dots, A_k)
$$

donde cada componente de $f$ es una concatenación de símbolos terminales y de componentes de componentes de las tuplas generadas por $A_1, \ldots, A_k$.

Para cada no terminal enriquecido:
$$
    A'_i  = A_i[(p_1^{(i)}, q_1^{(i)}), \ldots, (p_{d_i}^{(i)}, q_{d_i}^{(i)})]
$$

asociado a cada símbolo $A_i$ se añade la regla de producción  correspondiente
$$
    A'_0 \to f(A'_1, \dots, A'_k)
$$
en la nueva gramática, siempre que las transiciones del DFA sean compatibles con la concatenación definida por $f$.
Más precisamente, considérese una componente de $f$ de la forma:

$$
    \alpha_0 z_1 \alpha_1 z_2 \cdots z_m \alpha_m,
$$

donde cada $\alpha_j$ es una cadena de símbolos terminales y cada $z_j$ corresponde a una componente de algún argumento de $f$. La regla anterior se añade a la nueva gramática si:

La regla anterior se añade únicamente si:
\begin{itemize}
    \item el estado inicial asociado a $z_1$ se obtiene al procesar $\alpha_0$ desde el estado inicial de la componente correspondiente de $A_0$
    \item para cada \(i \ge 2\), el estado inicial asociado a \(z_i\) coincide con el estado alcanzado después de procesar \(z_{i-1}\alpha_{i-1}\)
    \item el estado final de la componente correspondiente de \(A_0\) se obtiene tras procesar \(z_m\alpha_m\).
\end{itemize}

En otras palabras, las reglas de la nueva gramática son exactamente aquellas derivaciones de la gramática original que corresponden a recorridos válidos del DFA.

% aqui quizas falte ponerle a la nueva regla la funcion identidad

Finalmente, por cada estado de aceptación $q_f$, se añade una nueva regla de la forma $S' \to S[(q_0, q_f)]$ que exige que la cadena generada lleve al DFA desde su estado inicial $q_0$ hasta un estado de aceptación $q_f$. De esta forma, la gramática construida genera exactamente el lenguaje $L(G) \cap R$, donde $G$ es la MCFG original y $R$ es el lenguaje reconocido por el DFA.

Como consecuencia, se tiene el siguiente teorema.

\begin{theorem}
    Sea $G$ una MCFG y $R$ un lenguaje regular. El lenguaje $L(G) \cap R$ es regular.
\end{theorem}

La construcción anterior incrementa el número de símbolos no terminales y reglas de producción al considerar todas las combinaciones posibles de estados del DFA asociadas a las componentes de cada no terminal.

Sea $q = |Q|$ el número de estados del DFA y sea $D = \max_{A \in N} dim(A)$ la máxima dimensión de un símbolo no terminal de la MCFG.

Si un símbolo no terminal $A$ tiene dimensión $d$, entonces los símbolos enriquecidos asociados a $A$ tienen la forma $A[(p_1, q_1), \ldots, (p_d, q_d)]$, donde cada $p_i, q_i \in Q$. Por lo tanto, existen $q^{2d}$ versiones enriquecidas posibles de $A$. En particular, el número de símbolos no terminales construidos para un símbolo de dimensión a lo sumo $D$ es $O(q^{2D})$, aunque, en la práctica, muchos de estos símbolos no se utilizan en reglas de la nueva gramática.

Considérese ahora una regla de producción $A_0 \to f(A_1, \ldots, A_k)$. Para construir las reglas correspondientes en la nueva gramática, es necesario considerar todas las combinaciones posibles de estados asociados a los símbolos $A_0, A_1, \ldots, A_k$. Si

$$
    \sum_{i=0}^{k} dim(A_i) \leq (k+1)D,
$$

entonces el número de combinaciones posibles de estados es
$$
    q^{2\sum_{i=0}^{k} dim(A_i)} \leq q^{2(k+1)D}.
$$

Sea $r = \max \{k \suchthat A_0 \to f(A_1, \ldots, A_k) \in P\}$, es decir, el número máximo de símbolos no terminales en el lado derecho de una regla de producción. Entonces, en el peor caso una regla original puede generar $O(q^{2(r+1)D})$ reglas en la nueva gramática.

Además, por cada combinación de estados, es necesario verificar la compatibilidad de las transiciones del DFA con la concatenación definida por la función de reordenamiento correspondiente. Sea $V$ el máximo número de variables que aparecen en una componente de una función de reordenamiento, y sea $T$ la máxima longitud de las cadenas de símbolos terminales insertadas entre variables. La verificación de la compatibilidad para una regla cuesta entonces $O(DVT)$, ya que a lo sumo deben verificarse $D$ componentes.

Por último, si $|P|$ es el número de reglas de producción en la MCFG original, el número total de reglas en la nueva gramática es $O(|P| \cdot q^{2(r+1)D})$, y el costo total de la construcción es $O(|P| \cdot q^{2(r+1)D} \cdot DVT)$.

% aquí me gustaría llegar a una conclusión, en plan, aunque es exponencial en dependencia de la gramática, es polinomial respecto al número de estados del DFA, lo cual es relevante para el capítulo 2, ya que el DFA se construye a partir de la fórmula. Escucho consejos :v

\subsection{Intersección de pMCFG con DFA}

Hallar la intersección de una pMCFG con un DFA es más complejo que en el caso de una MCFG, debido a que las funciones de reordenamiento no necesariamente son lineales. En \cite{Seki1991OnMC} se propone el mismo método que para una MCFG. Sin embargo, este no es suficiente, es posible demostrarlo con un contraejemplo.

Sea la pMCFG $G = (N, T, F, P, S)$ con $N = \{S, A\}, T = \{a\}, F = \{f, g\}$ y $P = \{S \to f(A), A \to g()\}$ donde $f(x) = xx$ y $g() = a$. Sea $\mathcal{A} = ({q_0, q_1, q_2}, T, \delta, q_0, {q_2})$ con función de transición $\delta$ dada por $\delta(q_0, a) = q_1$ y $\delta(q_1, a) = q_2$. Es fácil verificar que la cadena $aa$ pertenece a $L(G) \cap L(\mathcal{A})$, por tanto al efectuar el método propuesto para hallar la intersección, la gramática obtenida debería ser capaz de generar $aa$.

Al efectuar el método propuesto. La nueva gramática $G' = (N', T, F, P', S')$ tienes los símbolos no terminales $N' = \{S'\} \cup \{A[(p, q)], S[(p, q)] \suchthat p, q \in Q\}$

En las nuevas reglas de producción $P'$ las reglas correspondientes a $A \to g()$ son de la forma $A[(p, q)] \to g()$. Dado que $g() = a$, las reglas de esta forma solo serán válidas si $\delta(p, a) = q$, lo cual solo se cumple para $(q_0, q_1)$ y $(q_1, q_2)$. Por lo tanto, las reglas de esta forma en $P'$ son $A[(q_0, q_1)] \to g()$ y $A[(q_1, q_2)] \to g()$. Por tanto, solo los no terminales $A[(q_0, q_1)]$ y $A[(q_1, q_2)]$ pueden generar la cadena $a$.
% creo que debería escribir menos math inline...              ...tengo sueño Zzz...

Las reglas en $P'$ correspondientes a $S \to f(A)$ con $f(x) = xx$, son de la forma $S[(p, q)] \to f(A[(r, s)])$. Analizando paso a paso si la regla es válida dada la forma $a_0z_1a_1z_2a_2$ con $a_0 = a_1 = a_2 = \eword$ y $z_1 = z_2 = x$ de la única componente generada por $f$ se tiene:

\begin{itemize}
    \item El autómata debe cumplir que $\delta^*(p, a_0 = \eword) = r$, por lo tanto, dada la definición de $\delta$ se tiene que $p = r$.
    \item El autómata procesa la variable $z_1 = x$ suponiendo que $\delta^*(r, x) = s$.
    \item El autómata debe cumplir que $\delta^*(s, a_1 = \eword) = r$, por tanto dado la definición de $\delta$ se tiene que $s = r$.
    \item El autómata procesa la variable $z_2 = x$ suponiendo que $\delta^*(r, x) = s$.
    \item El autómata debe cumplir que $\delta^*(s, a_2 = \eword) = q$, por tanto dado la definición de $\delta$ se tiene que $s = q$.
\end{itemize}

Luego se tiene que $p = r = s = q$, por lo tanto, solo las reglas de la forma $S[(p, p)] \to f(A[(p, p)])$ pueden ser válidas. Sin embargo, dado que $A[(p, p)]$ no puede generar ninguna cadena terminal, se concluye que $S[(p, p)]$ tampoco puede generar ninguna cadena terminal. Por lo tanto, la gramática obtenida no puede generar la cadena $aa$, lo cual demuestra que el método propuesto no es suficiente para hallar la intersección de una pMCFG con un DFA.

% aqui va la demostración de que aún así las pMCFG son cerradas bajo intersección con lenguajes regulares. Coming soon...